\section{REVIEW PROCESS AND EVIDENCE BASE}
\label{sec:review_methods}
\label{sec:corpus_appraisal_results}

The technical synthesis connects reported results to the physical setting and
shared resources described in Section~\ref{sec:foundations_comparison}.
This section explains how we selected studies, extracted their results, and
assessed their support for comparison. Review reporting follows the Preferred
Reporting Items for Systematic Reviews and Meta-Analyses (PRISMA) 2020
statement and its search extension, PRISMA-S \cite{PRISMA2020,PRISMAS2021}.

\subsection{Search, Eligibility, and Study Reconciliation}
\label{sec:study_selection_reconciliation}

The primary evidence base covered communication and sensing within a shared
optical or photonic system. Eligible reports were peer-reviewed journal
articles or full conference papers with sufficient technical content in
English and assessable evidence for at least one survey domain. The
publication window extended from 1 January 2020 through 22 June 2026, using
the date when a report first became publicly available. Eligibility was
assessed separately from support for a joint quantitative comparison.
Contextual publications informed field positioning outside the primary study
count.

We searched Scopus, IEEE Xplore, ScienceDirect, SpringerLink, Wiley Online
Library, and Taylor \& Francis Online. The strategies covered O-ISAC phrases
and combinations of joint sensing--communication terms with optical or
photonic platform terms. All final searches were executed on 22 June 2026.
Electronic Supplement S-Search provides the available or reconstructed
strategies and export records. The query-to-export mapping for Taylor \&
Francis Online could not be reconstructed.

Selection proceeded from deduplicated records through title and abstract
screening to full-text assessment. The review was investigator supervised,
with artificial intelligence (AI) tools assisting record preparation,
extraction, and proposed decisions. The investigator authorized specified sets
of decisions, while rule-based checks examined source links and counting
units. Routine independent duplicate human review and third-reviewer
arbitration were not performed. Electronic Supplement S-Protocol describes the
executed workflow.

The review was retrospectively registered on the Open Science Framework on 12
February 2026 under record \texttt{7f6wb}. This registration documents an
earlier review state. Dated amendments in S-Protocol describe subsequent
changes to the search scope, report reconciliation, review conduct, and
synthesis.

Before synthesis, we linked reports describing the same underlying study to
avoid counting them as separate studies. This mapped 227 eligible reports to
206 studies, with 21 companion reports retaining their source links.
Figure~\ref{fig:prisma_report_study_flow} presents the record, report, and
study counts. Electronic Supplement ST-01 provides the report-to-study
mapping, and S-Exclusions records the eligibility decisions. These links
provide the basis for retaining the source and conditions of each extracted
result in Section~\ref{sec:extraction_analysis_units}.

\begin{figure*}[!t]
\centering
\includegraphics[width=\textwidth]{figures/fig_section3_prisma_layout.pdf}
\caption{Selection and reconciliation of the O-ISAC evidence base.
Six sources yielded 1,733 records, of which 1,259 entered title and abstract
screening. The retrieval route contained 330 unique reports after two
bibliographic aliases were consolidated. Full-text assessment covered
272 reports and produced 227 eligible reports, mapped to 206 studies.
The separate contextual corpus contains 67 records for field positioning.}
\label{fig:prisma_report_study_flow}
\end{figure*}

\subsection{Extraction and Units of Analysis}
\label{sec:extraction_analysis_units}

A study can contribute several results obtained under different conditions.
We therefore retained the report and source location for each extracted item.
Metric records preserved the reported quantity and unit, the measurement
point, and the test configuration. The configuration included operating
conditions, baselines, and validation settings where supplied by the source.
Electronic Supplement S-Data Dictionary defines the extraction fields.

Communication and sensing measurements remained separate unless the source
linked them to a shared setting. Relationship records captured the design
factors or objectives discussed in the source, together with their constraints
and reported effects. These records support the mechanism-based synthesis,
while Section~\ref{sec:joint_performance_interpretation} establishes the
conditions for interpreting a numerical joint operating point.
Table~\ref{tab:evidence_reconciliation} distinguishes these analysis units.

\begin{table*}[!t]
\caption{Units of analysis and their roles in the technical synthesis}
\label{tab:evidence_reconciliation}
\centering
\begingroup
\footnotesize
\setlength{\tabcolsep}{4pt}
\renewcommand{\arraystretch}{1.12}
\begin{tabularx}{\textwidth}{@{}
>{\raggedright\arraybackslash}p{0.17\textwidth}
>{\raggedleft\arraybackslash}p{0.06\textwidth}
>{\raggedright\arraybackslash}p{0.35\textwidth}
>{\raggedright\arraybackslash}X@{}}
\toprule
\textbf{Analysis unit} & \textbf{Count} &
\textbf{Meaning of one entry} & \textbf{Use in the survey} \\
\midrule
Included study & 206 &
An underlying study after reconciliation of its eligible reports. &
Study-level platform, integration, and validation profiles. \\
\addlinespace[3pt]
Primary coding record & 8,203 &
An extracted evidence item, metric record, or tradeoff record retained for
primary synthesis. &
Source-linked technical analysis, with each record retaining its own unit
and conditions. \\
\addlinespace[3pt]
Metric record & 4,779 &
A reported measure with its definition and evaluation setting. &
Communication, sensing, and implementation performance in
Section~\ref{sec:tradeoffs}. \\
\addlinespace[3pt]
Substantive relationship & 402 &
A source-described relation among design factors, objectives, or outcomes. &
Mechanisms and condition-dependent tradeoffs in
Section~\ref{sec:tradeoffs}. \\
\addlinespace[3pt]
Synthesis group & 115 &
A set of studies grouped for thematic synthesis. &
111 theme-specific groups support synthesis. Four broad classification
groups retain coverage information. \\
\bottomrule
\end{tabularx}
\par\vspace{3pt}
\noindent\parbox{\textwidth}{\raggedright\footnotesize
The counts refer to different analysis levels. Metric and relationship records
contribute to the primary coding total, while a study may contribute to several
synthesis groups. The primary total comprises 3,020 evidence items,
4,779 metric records, and 404 tradeoff records. Of the latter, 402 describe
relationships and two document the absence of a reported tradeoff.\par}
\endgroup
\end{table*}

Extraction retained numerical values stated by the source. Graphical results
supported interpretation in their reported form, without creating additional
numerical measurements. Source conflicts were handled at the individual-claim
level, so supported findings from the same study remained available.
The complete register contains 8,306 records, including 31 contextual items
and 72 items with source conflicts outside the primary synthesis.
Electronic Supplement S-Evidence provides the primary synthesis records and
their source links. S-Protocol describes the rules for handling contextual
items and source conflicts.

\subsection{Technical Appraisal and Synthesis}
\label{sec:technical_appraisal_synthesis}

We assessed all included studies using the Technical Quality Assessment
Framework (TQAF) developed for this review. Its deterministic rules score
eight technical and reporting dimensions from 0 to 3. Overall evidence
contribution is recorded separately. TQAF is a review-specific descriptive
appraisal whose scoring has not been independently validated. It characterizes
support for interpretation and reuse, rather than providing a conventional
risk-of-bias assessment.

Figure~\ref{fig:tqaf_profile} separates the appraisal dimensions by their
role in reading the evidence. Metric clarity and reporting completeness help
establish what a result means in its original setting. Validation and
reproducibility address how that result was tested and can be reconstructed.
Benchmark readiness and comparison admissibility address its use beyond the
original experiment. Electronic Supplements S-Protocol and S-Appraisal provide
the scoring rules and study-level records, respectively.

\begin{figure*}[!t]
\centering
\includegraphics[width=\textwidth]{figures/final_submission_2026-08-25/fig04_tqaf_profile.pdf}
\caption{Technical appraisal of the 206 included studies under the
review-specific TQAF rules. Each row groups all studies by strong, adequate,
or low support. The interpretation dimensions concern the meaning and
reporting of results in their source setting. Transfer and reuse dimensions
concern validation and reconstruction, together with benchmark and comparison
support. Overall evidence contribution is displayed separately.}
\label{fig:tqaf_profile}
\end{figure*}

The synthesis grouped studies by technical theme across seven survey domains.
Within each theme, we examined recurring mechanisms and reported performance
relations while preserving their physical and measurement conditions.
Thematic groups supported synthesis conclusions, while broad classification
groups recorded coverage across mixed topics. Electronic Supplement S-Bodies
provides group membership and the review-specific appraisal rules.

We used structured narrative synthesis to connect these heterogeneous
results. Numerical examples retain their original tasks, units, and operating
conditions. The review did not fit a pooled effect model or perform formal
publication-bias and sensitivity analyses. The resulting conclusions concern
supported mechanisms and bounded performance relations.

Section~\ref{sec:optical_modality_families} first examines the physical role
of optics and the resources shared by communication and sensing.
Those distinctions provide the context for the performance relations in
Section~\ref{sec:tradeoffs} and their validation in
Section~\ref{sec:validation_reproducibility}.
